\documentclass[a4paper,fleqn]{article} % added fleqn
\AtBeginDocument{\begin{multicols}{3}}
\AtEndDocument{\end{multicols}}
\usepackage[margin=0.2cm]{geometry}
\usepackage{amsmath,amssymb,amsfonts}
\usepackage{xcolor}
\usepackage{comment}
\usepackage{amsmath, amssymb}
\setlength{\mathindent}{0pt} % force no indent for displayed math
\usepackage{multicol}
\setlength{\columnsep}{0.1cm}

% reduce vertical space around displayed math and between gather lines
\setlength{\abovedisplayskip}{0pt}
\setlength{\abovedisplayshortskip}{0pt}
\setlength{\belowdisplayskip}{0pt}
\setlength{\belowdisplayshortskip}{0pt}
\setlength{\jot}{0pt}
\setlength{\parskip}{0pt}
\setlength{\parindent}{0pt}
\setlength{\parsep}{0pt}
\setlength{\partopsep}{0pt}
\setlength{\topsep}{0pt}
\setlength{\itemsep}{0pt}

\usepackage{etoolbox}
\makeatletter
\patchcmd{\section}{\vspace{1.5ex plus .2ex}}{\vspace{0.2ex}}{}{}
\patchcmd{\subsection}{\vspace{0.01ex plus .2ex}}{\vspace{0.1ex}}{}{}
\makeatother

\usepackage{titlesec}
\titlespacing*{\section}{0pt}{0ex}{0ex}
\titlespacing*{\subsection}{0pt}{0ex}{0ex}





\begin{document}

\section*{Teorija polj}
\subsection*{Vektorski operatorji}
\begin{gather*}
    \text{s = skalarno polje, V = vektorsko polje} \\
    \text{grad: } \nabla s = \left(\frac{\partial s}{\partial x}, \frac{\partial s}{\partial y}, \frac{\partial s}{\partial z}\right) \\
    \text{div: } \nabla \cdot \mathbf{V} = \frac{\partial V_x}{\partial x} + \frac{\partial V_y}{\partial y} + \frac{\partial V_z}{\partial z} \\
    \text{rot: } \nabla \times \mathbf{V} =
    \begin{vmatrix}
        \mathbf{i} & \mathbf{j} & \mathbf{k} \\
        \frac{\partial}{\partial x} & \frac{\partial}{\partial y} & \frac{\partial}{\partial z} \\
        V_x & V_y & V_z 
    \end{vmatrix} = \\
    = (Py-Qz,Pz-Rx,Qx-Py)\\
    \text{Laplacijan: } \nabla^2 s = \frac{\partial^2 s}{\partial x^2} + \frac{\partial^2 s}{\partial y^2} + \frac{\partial^2 s}{\partial z^2} \\
    \text{vek. produkt treh vektorjev: BAC-CAB}\\ \nabla \times (\nabla \times \vec{v}) = \nabla (\nabla \cdot \vec{v}) - \vec{v}(\nabla \cdot \nabla)) \\
    \text{Identiteta dvokratnega rotorja: }\\
    rot(rot \vec{v}) = grad(div \vec{v}) - \nabla^2 \vec{v} 
\end{gather*}

\begin{comment}
    \subsection*{Lastnosti polj}
    \begin{gather*}
        \text{Irotacijalno: }  \nabla \times \mathbf{F} = 0 \\
        \text{Solenoidno: } \nabla \cdot \mathbf{F} = 0 \\
        \text{Potencial: }  \vec{v}(x,y,z) = grad u(x,y,z) \\
        \text{u = skalarno p. }  \text{in potencial polja v}
    \end{gather*}
\end{comment}


\section*{Prostorska geometrija}
\subsection*{Koordinatni sistemi}
\begin{gather*}
    \text{Polarni: } x = r\cos\theta, \quad y = r\sin\theta \quad J = r\\
    \text{Cilindrični: }\\ x = r\cos\theta, \quad y = r\sin\theta, \quad z = z \quad J = r\\
    \text{Sferični: } x = r\sin\phi\cos\theta, \quad y = r\sin\phi\sin\theta,\\
    \quad z = r\cos\phi \quad J = r^2\sin\phi \\
    J = det \begin{bmatrix}
            \frac{\partial x}{\partial u} & \frac{\partial x}{\partial v} \\
            \frac{\partial y}{\partial u} & \frac{\partial y}{\partial v} 
            \end{bmatrix} 
    = det \begin{bmatrix}
            x _u & x _v \\
            y _u & y _v 
            \end{bmatrix} 
\end{gather*}

\subsection*{Razno formule}
\begin{gather*}
    \text{Ločna dolžina: } s = \int_a^b |\mathbf{r}'(t)| \, dt=\\ = \int_a^b \sqrt{x'(t)^2 + y'(t)^2 + z'(t)^2} \, dt \\
    \text{parame. točk: } \vec{r} = r _A * t(r _B - r _ A)\\
    \text{Odvod integrala s parametrom: } \\
    F(x)=\int_{u(x)}^{v(x)} f(x,y)dy \quad \Rightarrow \\
    F'(x)=\int_{u(x)}^{v(x)} \frac{\partial(f(x,y))}{\partial(x)}dy+\\+f(x,v(x))*v'(x)-f(x,u(x))*u'(x) \\
    \text{Naravna parameterizacija: } \\ |\mathbf{r}'(s)|^2 = 1, \quad s = \text{ločna dolžina} \\
    \text{Tangentna premica in normalna ravnina: } \\ \vec{e}\text{ ali }\vec{n} = \dot\vec{r}(t0) \\
    \text{Kot med Vektorjema: } \cos\theta = \frac{\vec{a}\cdot\vec{b}}{|\vec{a}||\vec{b}|}
\end{gather*}

\subsection*{Večkratni integrali}
\begin{gather*}
    \text{Dvojni integral: } \\ \iint_D f(x,y) \, dA = \int_a^b \int_{g_1(x)}^{g_2(x)} f(x,y) \, dy \, dx\\
    \text{Ploščina: } A = \iint_D 1 \, dA\\
    \text{Prostornina: } V = \iint_D f(x,y) \, dA\\
    \text{Težišče: } \bar{x} = \frac{1}{A}\iint_D x \, dA, \quad \bar{y} = \frac{1}{A}\iint_D y \, dA\\
    \text{Masa: } m = \iiint_E \rho(x,y,z) \, dV\\
    \text{Vztrajnostni moment: } I = \iiint_E r^2 \rho \, dV
\end{gather*}

\subsection*{Krivuljni integrali}
\begin{gather*}
    \textbf{\scriptsize S. polje:} 
    \int_C f(x,y,z) ds =\\=\int_{t0}^{t1} f(x(t),y(t),z(t))*\sqrt{\dot{x}^2,+\dot{y}^2+\dot{z}^2} dt \\
    \textbf{\scriptsize V. polje:}
    \int_C \mathbf{V} \cdot d\mathbf{r} = \int_C(Pdx + Qdy + Rdz)= \\ =\int_{t0}^{t1}(P\dot{x}+Q\dot{y}+R\dot{z}) dt \\
    \quad d\vec{r} = \dot\vec{r} dt = (\dot{x}dt,\dot{y}dt,\dot{z}dt)\\
    \text{Neodvisen od poti če }\vec{v}\text{ potencialno polje } \\ (rot\vec{V}=0) 
\end{gather*}

\subsection*{Ploskovni integrali}
\begin{gather*}
    \textbf{\scriptsize S. polje:}
        \iint_S f \, dS =\\= \iint_D f(\mathbf{r}(u,v)) \sqrt{EG-F^2} \, du \, dv \\
        \text{kjer } E = \mathbf{r}_u \cdot \mathbf{r}_u,\quad F = \mathbf{r}_u \cdot \mathbf{r}_v,\quad G = \mathbf{r}_v \cdot \mathbf{r}_v \\
            z = z(x,y), \quad \vec{n} = (-z _x, -z _y, 1)\\
    \textbf{\scriptsize V. polje:}
        \iint_S \mathbf{V} \cdot d\mathbf{S} =\\= \iint _S \mathbf V \cdot \vec{v} \text{ } dS = \iint_D [\mathbf{V}, \mathbf{r}_u, \mathbf{r}_v] \, du \, dv \\
            \vec{v} = \frac{\vec r _u \times \vec r _v}{|\vec r _u \times \vec r _v|}, \quad
            |\vec r _u \times \vec r _v| = \vec n
\end{gather*}

\subsection*{Integralski izreki}
\begin{gather*}
    \textcolor{red}{Greenov, Stokesov, Gaussov: }\\
     \oint_C (P \, dx + Q \, dy) = \iint_D \left(\frac{\partial Q}{\partial x} - \frac{\partial P}{\partial y}\right) dxdy\\ 
    \text{ \tiny C : zaključena pozitivno orientirana krivulja, ki je rob območja D}\\
    \oint_C \mathbf{V} \cdot d\mathbf{r} = \iint_S (\nabla \times \mathbf{V}) \cdot d\mathbf{S}\\
    \text{ \tiny C : rob površine S}\\
    \iint_S \mathbf{V} \cdot d\mathbf{S} = \iiint_E \nabla \cdot \mathbf{V} \, dV\\
    \text{ \tiny S : zaprt površina, ki je rob območja E}
\end{gather*}

\section*{Kompleksna analiza}

\subsection*{Elementarne komp. fun.}
\begin{gather*}
    z = x + iy; \quad |z| = \sqrt{x^2 + y^2}; \quad\\
     \arg z = \arctan\left(\frac{y}{x}\right)\\
    e^z = e^x(\cos y + i\sin y); \quad d \varphi = \frac{dz}{iz}\\
    \log z = \log|z| + i(\arg z + 2\pi k)\\
    \sin z = \frac{e^{iz} - e^{-iz}}{2i}, \cos z = \frac{e^{iz} + e^{-iz}}{2}\\
    \sinh z = \frac{e^z - e^{-z}}{2}, \cosh z = \frac{e^z + e^{-z}}{2}
\end{gather*}

\subsection*{Integral Kompleksne funkcije}
\begin{gather*}
    \int_C f(z) \, dz = \int_C (u+iv)(dx+idy) = \\= \int_C (u \, dx - v \, dy) + i \int_C (v \, dx + u \, dy)\\
    = \int_{t_a}^{t_b} f(z(t)) \cdot \dot{z}(t) \, dt\\
    \text{Cauchjeva form.: } f(z_0) = \frac{1}{2\pi i} \oint_C \frac{f(z)}{z - z_0} dz\\
    \text{!Integral analitične. fun. ni odvisen od poti!}
\end{gather*}

\subsection*{Taylorjeva in Laurentova vrsta}
\begin{gather*}
        f(x) = \sum_{n = 1}^{\infty} \frac{f^{(n)}(x_0)}{n!} (x-x_0)^n\\
        \text{Geometrijska, sin, cos, e vrsta: }\\
        \frac{1}{1-z}=\sum_{n=0}^{\infty}z^n \quad ; \text{$|z| < 1$}\\
        \sin z = \sum_{n=0}^{\infty}(-1)^n\frac{z^{2n+1}}{(2n+1)!} \quad ; \text{za vse z}\\
        \cos z = \sum_{n=0}^{\infty}(-1)^n\frac{z^{2n}}{(2n)!} \quad ; \text{za vse z}\\
        e^z = \sum_{n=0}^{\infty}\frac{z^n}{n!} \quad ; \text{za vse z}\\
        \text{Laurentova vrsta: }\\
        f(z) = \sum_{n=-\infty}^{\infty} a_n (z - z_0)^n \\ a_n = \frac{1}{2\pi i} \oint_C \frac{f(z)}{(z - z_0)^{n+1}} dz
\end{gather*}

\subsection*{Analitičnost in odvedljivost in harmoničnost}
\begin{gather*}
    \text{C-R pogoj: }\\
    \frac{\partial u}{\partial x} = \frac{\partial v}{\partial y}, \quad \frac{\partial u}{\partial y} = -\frac{\partial v}{\partial x}\\
    u = \int v _y \, dx = - \int v _x \, dy + C\\
    \text{Pogoj harmoničnosti: }\\ 
    \nabla^2 u = \nabla^2 v = 0 // w _{xx} + w _{yy} = 0
\end{gather*}

\subsection*{Formule kompleksne analize}
\begin{gather*}
    \text{Residum: } \oint_C f(z) dz = 2\pi i \sum \text{Res}(f, z_k)\\
    Res f(z) = \frac{1}{(m-1)!} *\\* \lim_{z \to z_0}[ \frac{d^{m-1}}{dz^{m-1}} \left((z - z_0)^m f(z)\right)] \\
    \text{[m]-pol n-te stopnje}\\
    \text{L.V.:} f(z) = \sum_{n=-\infty}^{\infty} C_n (z - z_0)^n \\
    C_{-1} = Res f(z) \, || \, f'(z) \text{ za stopnje ničel} \\
    \text{Eul. form. v } \mathbb{C}: \text{(realna rac. funkcija [f])} \\
    sin(\phi) = \frac{e^{i\phi} - e^{-i\phi}}{2i} \\
    cos(\phi) = \frac{e^{i\phi} + e^{-i\phi}}{2} \\
    z = e^{i\phi}|d\phi = {dz}/{iz}| \int_{0}^{2\pi}f(sin(\phi),cos(\phi))d\phi\\
    \frac{z - z^{-1}}{2i} = sin(\phi)\, | \,\frac{z + z^{-1}}{2} = cos(\phi)
\end{gather*}

\subsection*{Konformne presl. (C-R), Lomljena Linearna presl. / Mobiusova presl.}
\begin{gather*}
\text{Preslikava: } w = f(z) = u(x,y) + iv(x,y)\\
f(z) = \frac{az + b}{cz + d} ; ad - bc \neq 0\\
\end{gather*}

\newpage

\section*{\small Integrali}

\subsection*{\footnotesize Pravila integriranja}
\begin{align*}
\text{Substitucija: } &\int f(g(x))g'(x) dx = \int f(u) du, \quad u = g(x)\\
\text{Po delih: } &\int u \, dv = uv - \int v \, du\\
\text{Parcialni ulomki: } &\frac{P(x)}{Q(x)} = \sum \frac{A_i}{(x-a_i)^{n_i}}
\end{align*}

\subsection*{\footnotesize Elementarni integrali}
\begin{align*}
\int x^n dx &= \frac{x^{n+1}}{n+1} + C, \quad n \neq -1\\
\int \frac{1}{x} dx &= \ln|x| + C\\
\int e^x dx &= e^x + C\\
\int \sin x dx &= -\cos x + C\\
\int \cos x dx &= \sin x + C\\
\int \frac{1}{1+x^2} dx &= \arctan x + C\\
\int \frac{1}{\sqrt{1-x^2}} dx &= \arcsin x + C
\end{align*}

\section*{\small Odvodi}

\subsection*{\footnotesize Pravila odvajanja}
\begin{align*}
(f + g)' &= f' + g'\\
(fg)' &= f'g + fg'\\
\left(\frac{f}{g}\right)' &= \frac{f'g - fg'}{g^2}\\
(f(g(x)))' &= f'(g(x)) \cdot g'(x)
\end{align*}

\subsection*{\footnotesize Elementarni odvodi}
\begin{align*}
\frac{d}{dx}x^n &= nx^{n-1}\\
\frac{d}{dx}e^x &= e^x\\
\frac{d}{dx}\ln x &= \frac{1}{x}\\
\frac{d}{dx}\sin x &= \cos x\\
\frac{d}{dx}\cos x &= -\sin x\\
\frac{d}{dx}\tan x &= \sec^2 x\\
\frac{d}{dx}\arctan x &= \frac{1}{1+x^2}\\
\frac{d}{dx}\arcsin x &= \frac{1}{\sqrt{1-x^2}}
\end{align*}

\subsection*{Trigonometrske funkcije}
\begin{gather*}
    \sin^2 x + \cos^2 x = 1\\ 
    \tan x = \frac{\sin x}{\cos x}\\
    \sin 2x = 2\sin x \cos x\\
    \cos 2x = \cos^2 x - \sin^2 x \\
    \tan 2x = \frac{2\tan x}{1 - \tan^2 x}\\
    \sin(x \pm y) = \sin x \cos y \pm \cos x \sin y\\
    \cos(x \pm y) = \cos x \cos y \mp \sin x \sin y\\
    \tan(x \pm y) = \frac{\tan x \pm \tan y}{1 \mp \tan x \tan y}\\
    \text{Adicijski izrek za sin,cos}: \\
    \sin (x+iy)=\sin x \cosh y + i \cos x \sinh y\\
    \cos (x+iy)=\cos x \cosh y - i \sin x \sinh y
\end{gather*}

\end{document}